\documentclass[10pt]{article}
\usepackage[margin=.88in]{geometry}
\usepackage{amsmath,amssymb}
\usepackage[colorlinks=true,urlcolor=blue,linkcolor=black]{hyperref}
\setlength{\parskip}{.45em}
\setlength{\parindent}{0pt}
\setlength{\emergencystretch}{3em}
\providecommand{\tightlist}{}
\renewcommand{\labelitemi}{$\bullet$}

\setlength{\emergencystretch}{3em}
\title{Fixed-deadline coupon coverage:\\balanced draws and late deadlines}
\author{Leo Gao\\\small AI-assisted research}
\date{}
\begin{document}
\pdfmapfile{+cm.map}\pdfmapfile{+symbols.map}
\maketitle

\section*{Fixed-deadline coupon coverage for all balanced package
laws}\label{fixed-deadline-coupon-coverage-for-all-balanced-package-laws}

\textbf{Theorem.} Let \(n=qs\), with positive integers \(q,s\). At each
draw choose a package of at most \(s\) coupons. Suppose, conditional on
every past history, each coupon has inclusion probability \(1/q\). Then
at every integer deadline \(m\), the probability of full coverage is at
most \(q!{m\brace q}/q^m\). The iid uniform law on the blocks of an
\(s\)-set partition attains the bound. Among iid balanced laws it is the
unique maximizer up to relabeling at each feasible deadline \(m\ge q\).
For iid draws this proves the full \emph{balanced-marginal case} of
\href{https://hexagonmath.org/pdf/2611.00814v1}{Long's Conjecture 8.1}.
Unequal marginals are not treated here. Each draw necessarily has size
\(s\) since its expected size meets the cap.

\textbf{Proof strategy.} The normalized count \(Y\) of covered coupons
increases by at most one with conditional mean increment \((q-Y)/q\).
The partition law has the same drift but integer jumps. We construct a
family of convex functions for a one-step comparison. Linear
interpolation at integer states does not suffice. The functions below
interpolate hypergeometric probabilities by starting a
without-replacement recurrence from a fractional first step.

Write \((v)_r=v(v-1)\cdots(v-r+1)\), \((v)_0=1\). Put \(b_0=1\) and, on
integers \(0\le j\le q\), put \(b_r(j)=(j)_r/(q)_r\) for
\(1\le r\le q\). On \(y=j+u\), \(0<u<1\), \(0\le j\le q-1\), set \[
 b_r(j+u)=\frac{(j)_{r-1}}{(q-u)_r}
      \left[j-r+1+\frac{ru}{q}(q-u-j)\right].                 \tag{1}
\] In particular \(b_1(y)=y/q\). For \(j<r-1\), the expression is zero.
For \(j\ge r-1\), it is nonnegative. At \(u=0\) the integer value is
immediate. If \(r<q\), at \(u=1\) the bracket becomes \((j+1)(q-r)/q\),
giving the next integer value; for \(r=q\) expression (1) vanishes up to
\(j=q-1\) and subsequently reduces to a positive constant times
\(u/(q-u)_{q-1}\), tending to one at \(y=q\).

\subsection*{The recurrence before the last unit
interval}\label{the-recurrence-before-the-last-unit-interval}

For \(0\le y\le q-1\), \[
 (q-y)\{b_r(y+1)-b_r(y)\}=r\{b_{r-1}(y)-b_r(y)\}.           \tag{2}
\] Here is a derivation that explains (1). For fixed \(u\in[0,1)\)
define \[
 c_i(j)=(1-iu/q)\,(q-u-i)_j/(q-u)_j,\qquad
 B_r(j)=\sum_{i=0}^{r}(-1)^i {r\choose i}c_i(j).
\] These \(c_i\) satisfy \((q-u-j)[c_i(j+1)-c_i(j)]=-i c_i(j)\), and
\(i{r\choose i}=r({r\choose i}-{r-1\choose i})\), with zero-extended
binomial coefficients. This gives (2) for the \(B_r\)'s. To see
\(B_r(j)=b_r(j+u)\) (at the stated arguments), write \(Q=q-u\). The
identity \[
 \sum_{i=0}^r(-1)^i{r\choose i}\frac{(Q-i)_j}{(Q)_j}
  =\frac{(j)_r}{(Q)_r}
\] is inclusion-exclusion for hitting \(r\) specified points in a
\(j\)-element sample without replacement when \(Q\) is a large integer;
it continues as a rational function. The \(i\)-weighted sum divided by
\(r\) is the difference of the unweighted sums for \(r\) and \(r-1\).
Hence \(B_r=(1-ru/q)(j)_r/(Q)_r+(ru/q)(j)_{r-1}/(Q)_{r-1}\), exactly
(1). The integer endpoint \(y=q-1\) follows directly from the integer
values.

\subsection*{Convexity}\label{convexity}

We prove \(b_r\) is convex. Only \(r\ge2\) needs attention. Within a
unit interval with \(j\ge r-1\), drop the positive factor \((j)_{r-1}\)
in (1). The resulting expression is affine in the integer parameter
\(j\), so it is enough to prove interval convexity at \(j=r-1\) and
\(j=q-1\). For \(j=r-1\) it is \((ru/q)/(q-u)_{r-1}\), a product of
nonnegative, nondecreasing, convex functions. This also covers \(r=q\).

For \(j=q-1\), set \(t=q-r\ge1\) and \(D=\prod_{v=t+1}^{q}(v-u)^{-1}\).
The remaining function is \(ND\), \(N=t+(r/q)u(1-u)\). Write
\(S_k=\sum_{v=t+1}^q(v-u)^{-k}\). Differentiation gives \[
 (ND)''/D=N(S_1^2+S_2)+2(r/q)(1-2u)S_1-2r/q.              \tag{3}
\] For \(t\ge2\), set \(z=t+(r+1)/2-u\), the arithmetic mean of the
\(v-u\). Then \(S_1\ge r/z\) and \(S_2\ge S_1^2/r\). In the resulting
lower bound \((1+1/r)NS_1^2+2(r/q)(1-2u)S_1-2r/q\), the quadratic is
increasing for \(S_1\ge r/q\): indeed \(N\ge t\ge1\). Substituting
\(r/z\) and multiplying by \(qz^2/r>0\) leaves \[
 G(u)=(r+1)[qt+ru(1-u)]+2r(1-2u)z-2z^2.
\] Its quadratic coefficient is \(-(r-1)(r-2)\le0\). Its endpoint
values, after dividing by \((r-1)/2\), are \(2rt+r+2t^2+4t+1\) and
\(2rt-3r+2t^2-4t+1\), both positive at \(t\ge2\). This proves (3) there.

For \(t=1\), multiply (3) by \(q/r\). Write \(d=\min(r,5)\),
\(L_k=\sum_{v=2}^{d+1}(v-u)^{-k}\), and let \(c=1+1/r\) for \(r<5\),
\(c=1\) otherwise. A lower bound is \[
 [c+u(1-u)](L_1^2+L_2)+2(1-2u)L_1-2.                       \tag{3a}
\] This is equality for \(r<5\). For \(r\ge5\), first reduce the
coefficient to \(c+u(1-u)\), then replace \(S_2\) by \(L_2\) and \(S_1\)
by \(L_1\); the latter decreases the expression since
\(L_1\ge\sum_{v=2}^6 1/v>1\).

For completeness, (3a) is positive by the following explicit polynomial
certificate. Substitute \(u=x/(1+x)\), \(x\ge0\), and multiply by
\(a_d\prod_{i=1}^d(ix+i+1)^2\), where \((a_2,a_3,a_4,a_5)=(1,3,2,1)\).
In increasing degree the resulting coefficients are:

\begin{itemize}
\tightlist
\item
  \(d=2\): 45, 99, 72, 19, 1.
\item
  \(d=3\): 3968, 14880, 22200, 16702, 6540, 1200, 68.
\item
  \(d=4\): 92300, 497350, 1143965, 1462904, 1131922, 537956, 150953,
  22138, 1196.
\item
  \(d=5\): 1811232, 12643200, 38944680, 69530280, 79352100, 60087366,
  30229390, 9768970, 1853430, 162620, 1592.
\end{itemize}

All are positive; the case \(u=1\) also follows by continuity. They can
be reproduced without symbolic algebra by multiplying the denominators,
or by the independent exact-evaluation verifier. The larger 15,232-cell
rational enclosure in the supplement is an additional check, not
required by the proof.

At an integer \(j\) with \(r\le j\le q-1\) (\(r<q\)), direct
differentiation of (1) at its two ends gives the slope jump \[
 b_r'(j+)-b_r'(j-)= b_r(j)\,
 \frac{r(q-j)(r-1)(q-j-1)}{q(q-r)j(j-r+1)}\ \ge0.
\] For example, dividing the jump by \(r b_r(j)\) gives
\(-1/[q(q-r)]+(q-j)/[q(j-r+1)]-(q-j-1)/[j(q-r)]\), yielding the
displayed factor. At \(j=r-1\) the left slope is zero and the right
slope is nonnegative; earlier intervals are zero. This completes global
convexity.

\subsection*{The shortened last jump}\label{the-shortened-last-jump}

When \(q-1<y<q\), the largest possible increment is \(q-y\), not one. We
need \[
 1+(r-1)b_r(y)\ \le\ r b_{r-1}(y).                           \tag{4}
\] For \(r=1\) this is equality. Otherwise set \(j=q-1\), \(w=1-u\),
\(t=q-r\), and \(P(w)=(t+w)\prod_{v=t+1}^{q-1}(1+w/v)\). Using (1),
multiplying (4) by \(P(w)>0\) and collecting terms gives \[
 P(w)\ \le\ t+rw-\frac{r(r-1)}{q(t+1)}w(1-w)^2.
\] Indeed \(P(w)b_r=t+(r/q)w(1-w)\) and
\(P(w)b_{r-1}=(t+w)[t+1+((r-1)/q)w(1-w)]/(t+1)\). The polynomial \(P\)
has nonnegative coefficients, \(P(0)=t,\ P(1)=q\). Its deficit below
this endpoint chord is at least \(w(1-w)\) times the sum of its
coefficients of degrees at least two, since \(w-w^k\ge w(1-w)\) for
\(k\ge2\). The coefficient sum is \[
 r-P'(0)=\sum_{i=1}^{r-1} i/(t+i).
\] If \(t=0\) this is \(r-1=r(r-1)/[q(t+1)]\). If \(t\ge1\) it is at
least \(r(r-1)/[2(q-1)]\ge r(r-1)/[q(t+1)]\). This proves (4).

\subsection*{Applying the kernel}\label{applying-the-kernel}

For \(0\le Y=y\le q-1\), convexity bounds \(b_r(Y')\) above by its chord
on \([y,y+1]\), where \(Y'\) is the state after one more draw. Taking
conditional expectation and applying (2) yields \[
 \mathbb E[b_r(Y')\mid\text{history}]
 \le (1-r/q)b_r(y)+(r/q)b_{r-1}(y).                         \tag{5}
\] For \(q-1<y<q\) use the chord on \([y,q]\) with interpolation weight
\((\mathbb E Y'-y)/(q-y)=1/q\), and then (4); at \(y=q\) both sides are
one. The required mean increment follows by summing the inclusion
probabilities over missing coupons. No Markov or independence assumption
is needed beyond the stated conditional marginals.

Start at \(Y_0=0\). For \(r\ge1\), \(b_r(0)=0\). The probability
\(v_r(m)\) that \(r\) specified labels have all appeared in \(m\)
ordinary iid uniform draws from \(q\) labels has these same initial
values and obeys (5) with equality: condition on whether the last draw
is one of the specified labels. Induction gives
\(\mathbb E b_r(Y_m)\le v_r(m)\) for every \(r\). Since \(b_q\ge0\) and
\(b_q(q)=1\), the success indicator is bounded by \(b_q(Y_m)\). Taking
\(r=q\) proves the theorem. The ordinary surjection count gives the
stated value. For \(q=1\) full coverage is forced on the first draw.

\subsection*{Strictness for iid laws}\label{strictness-for-iid-laws}

Suppose \(q\ge3\). For the first two packages \(A,B\), put
\(\xi=1-|A\cap B|/s\), so \(Y_2=1+\xi\). On that unit interval,
\(b_2(1+\xi)=2\xi/[q(q-\xi)]\). The chord gap in (5) at step two is
therefore \[
 \Delta=\frac{2}{q(q-1)}\,\mathbb E\!\left[\frac{\xi(1-\xi)}{q-\xi}\right].
\] If \(A,B\) can partially overlap, then \(\Delta>0\). Propagating this
deficit through (5) implies \[
 v_q(m)-\Pr(T\le m)\ \ge\ \Delta\Pr(\text{exactly two labels missing after }m-2
             \text{ ordinary uniform draws}).
\] Indeed the coefficient on the step-two deficit at \(r=2\) in the
iterated recursion is the transition probability of the chain
\(r\mapsto r-1\) with probability \(r/q\) (otherwise \(r\mapsto r\)),
started at \(q\) and run for \(m-2\) steps. It is positive for every
\(m\ge q\) since \(q-2\ge1\). If \(q=2\), any two supported partially
overlapping packages \(A,B\) give a strict final chord loss instead:
with positive probability the first \(m-1\) draws all equal \(A\) and
the final draw is \(B\); the last state is strictly inside \([1,2]\),
where \(b_2\) is strictly convex. Thus at every feasible \(m\ge2\),
partial overlaps are again impossible at equality. If no two supported
packages partially overlap, the support is a disjoint partition and
balance forces all its block probabilities to be \(1/q\). Conversely
those probabilities attain equality.

\textbf{Scope and verification.} This balanced result has no restriction
on the deadline and strengthens the class of processes to allow
conditional balance. It does \emph{not} say that an arbitrary optimizer
over unequal marginals must be balanced. For those iid laws see the
distinct large-deadline theorem in this bundle. The shorter three-block
proof gives a separate, elementary strictness term. Primary problem
statement: C. D. Long,
\href{https://hexagonmath.org/2611.00814v1}{\emph{Optimal coupon
packages and stochastic improvement of uniform group drawings}},
Conjecture 8.1; targeted searches found no prior all-deadline balanced
theorem, but no claim of priority or independent peer review is made.
\clearpage

\section*{The partition law uniquely maximizes late discrete coupon
coverage}\label{the-partition-law-uniquely-maximizes-late-discrete-coupon-coverage}

\textbf{Theorem.} Write \(n=qs\), \(q\ge2, s\ge1\). Independently draw
\(m\) random subsets of \([n]\) with common law \(\mu\), each of
cardinality at most \(s\) (the empty set is allowed). If \[
m\ge256n\log(2n),
\] the probability of covering \([n]\) is uniquely maximized, up to
relabeling, by drawing uniformly from \(q\) disjoint \(s\)-sets. The
maximum is \(q!\left\{\genfrac{}{}{0pt}{}{m}{q}\right\}/q^m\). The
logarithm is natural. For \(q=1\), the unique maximizing law for finite
\(m\ge1\) puts all mass on \([n]\): repeating any supported proper set
would fail.

The quantifiers are uniform: the competing law can change with the
deadline. This is a partial answer to Long's Conjecture 8.1, not a proof
at all deadlines. No computational assumption enters the proof.

\subsection*{1. A candidate maximizer has well-separated incidence
clusters}\label{a-candidate-maximizer-has-well-separated-incidence-clusters}

Let \(p=1/q, a=1-p\); \(E_i\) is the event coupon \(i\) is missing after
\(m\) draws, \(H=\Pr(\bigcup_i E_i)\), and \(p_i=\Pr(i\in A)\) for one
draw. For the partition law \(D\), \(H_D\le qa^m\) by the union bound.
Assume \(H\le H_D\); we prove \(\mu=D\) up to labels. Put \[
h_i=(1-p_i)/a,\qquad z_i=h_i^m.
\] Then \(z_i\le q\) and \(\sum_i h_i\ge n\), since \(\sum p_i\le s\).
Also \(p_i\ge1-aq^{1/m}\ge p-2\log(q)/m\ge p/2\): use \(e^u-1\le2u\) for
\(0\le u\le1\), and the stipulated bound on \(m\).

Join two coupons if the probability \(d_{ij}\) that \textbf{exactly one}
belongs to a draw is less than \(\delta=1/(4n)\). Every component has at
most \(s\) vertices. Otherwise take a connected tree on \(s+1\) of them
(grow a tree by one vertex at a time). With probability at least
\(p/2-s\delta>0\), its root is drawn and none of its \(s\) edges
disagrees; this would put \(s+1\) coupons in a package.

For coupons in different components, \(d_{ij}\ge\delta\). The
probability a draw misses both is \[
1-(p_i+p_j+d_{ij})/2\le a q^{1/m}-\delta/2.
\] Using \(a q^{1/m}\le2\) and \(1-u\le e^{-u}\) gives \[
\Pr(E_i\cap E_j)/a^m\le q\exp[-m/(16n)].
\] Let \(U_C=\bigcup_{i\in C}E_i\), and
\(R=n^2q\exp[-m/(16n)]\le n^3(2n)^{-16}\). Bonferroni on the components,
followed by the union bound for their pair intersections, gives \[
q\le s^{-1}\sum_i z_i\le\sum_C\Pr(U_C)/a^m\le H/a^m+R\le q+R. \tag{1}
\] The first step is Jensen and \(\sum h_i\ge n\); the second uses
\(|C|\le s\). Both the accumulated Jensen excess and the accumulated
within-component loss are therefore small.

In fact \(z_i>1/4\) for all \(i\). For \(u\ge0\),
\(g(u)=u^m-1-m(u-1)\ge0\); by (1), \(\sum_i g(h_i)\le sR<1/4\). On
\(0\le u\le4^{-1/m}\), \(g\) is decreasing and
\(g(4^{-1/m})=-3/4+m(1-4^{-1/m})\ge1/4\) for \(m\ge2\). The last bound
follows because \(v(1-e^{-b/v})\) is increasing for \(v>0, b=\log 4\),
and equality in this numerical bound holds at \(v=2\).

Each component has size exactly \(s\): otherwise its contribution
\(\Pr(U_C)/a^m-s^{-1}\sum_{i\in C}z_i\) exceeds \(1/(4s)>R\), using the
largest \(z_i\) in the component. Thus the components themselves give a
partition into \(q\) blocks.

\subsection*{2. Almost every package must be a whole
block}\label{almost-every-package-must-be-a-whole-block}

From (1) and nonnegativity of all component losses, for \(i,j\in C\) \[
\Pr(E_i\setminus E_j)/a^m\le [\Pr(U_C)-\Pr(E_j)]/a^m\le sR.
\] Write \(b=\Pr(j\in A,\ i\notin A)\). Independence gives \[
\Pr(E_i\setminus E_j)/a^m
=z_i\left[1-\left(1-\frac b{1-p_i}\right)^m\right]
\ge\tfrac14(1-e^{-mb}).
\] Thus \(mb\le-\log(1-4sR)\le8sR\) since \(4sR<1/2\). This bounds
\emph{each directed} within-block disagreement. The size deficiency also
satisfies \[
m(s-\sum p_i)/a\le \sum_i(z_i-1)\le sR
\] by the tangent inequality \(h_i^m\ge1+m(h_i-1)\).

Write \(\varepsilon\) for the probability of a package not equal to one
of the blocks. Such a package either has size less than \(s\), or splits
at least one block. Since \(s-|A|\) is integer, a union bound over
directed pairs yields \[
m\varepsilon\le asR+8n^2sR\le9 n^3R<1/(64q^2). \tag{2}
\] For the final numerical inequality use \(R\le2^{-16}n^{-13}\) and
\(2\le q\le n\).

\subsection*{3. A single defective draw is strictly
costly}\label{a-single-defective-draw-is-strictly-costly}

Conditioning on drawing a whole block gives a \(q\)-label law
\(\nu=(r_1,\ldots,r_q)\). Its least weight \(r=\min_j r_j\) is at least
\(1/(4q)\): a coupon in a block has
\(p_i\le(1-\varepsilon)r_j+\varepsilon\), whereas \(p_i\ge1/(2q)\) and
\(\varepsilon<1/(64q^2)\). Write \(F(t)\) for the completion CDF of this
single-label law, \(H_\nu(t)=1-F(t)\), and \(D_m=F(m)-F(m-1)\). By
insisting that a least-weight label be the last to arrive and
subtracting the chance of any other missing label, \[
D_m\ge r(1-r)^{m-1}
[1-(q-1)e^{-r(m-1)}]\ge(r/2)(1-r)^{m-1}.
\] Indeed, conditional on avoiding that label, each other label has
weight at least \(r/(1-r)\ge r\); \(m\ge256n\log(2n)\) ensures the
bracket is at least \(1/2\). As \(H_\nu(m)\le q(1-r)^m\), we obtain
\(D_m\ge H_\nu(m)/(8q^2)>0\).

Let \(K\sim\mathrm{Bin}(m,\varepsilon)\) count defective draws. Given
\(K=0\), completion has probability \(F(m)\). Given \(K=1\), it has
probability exactly \(F(m-1)\): no single defective package contains a
whole block, while good draws either cover an entire block or miss it
entirely. Given \(K\ge2\), even crediting the process certain success
only gains \(H_\nu(m)\) over the baseline \(F(m)\). If
\(\varepsilon>0\), then \[
\frac{\Pr(K\ge2)}{\Pr(K=1)}
\le \frac{\binom m2\varepsilon^2}{m\varepsilon(1-\varepsilon)^{m-1}}
\le\frac{m\varepsilon}{2(1-m\varepsilon)}
<1/(8q^2)
\] by (2). Mixing the three cases shows \(1-H<F(m)\).

Finally \(F(m)\le1-H_D\), strictly unless all positive block
probabilities equal \(1/q\). For completeness, average any two unequal
positive block weights while preserving their sum. Conditional on the
positions drawn from this pair and on the other labels, the probability
of seeing both in \(k\ge2\) pair positions is \(1-x^k-(1-x)^k\),
uniquely maximized at the equal split. For \(m\ge q\) and positive
weights, the case where every other label is present and \(k\ge2\) has
positive probability. The continuous function on the probability simplex
achieves a maximum; its boundary gives zero, and unequal interior
weights can be strictly improved, proving the assertion. Thus positive
\(\varepsilon\) contradicts \(H\le H_D\); at \(\varepsilon=0\) the same
strictness forces uniform block weights. This proves the theorem.

\subsection*{Source, scope,
obstruction}\label{source-scope-obstruction}

Christopher D. Long,
\href{https://hexagonmath.org/2611.00814v1}{\emph{Optimal coupon
packages and stochastic improvement of uniform group drawings}}, Hexagon
preprint v1, Nov 21 2026,
\href{https://hexagonmath.org/pdf/2611.00814v1}{Conjecture 8.1 and
Proposition 3.4}. Long explicitly asks for every discrete deadline,
proves the first feasible one \(m=q\), and obtains stochastic optimality
after Poissonization. His incidence-class argument for equality of the
continuous-time optimum is relevant prior work. Neither a transform
inequality nor dominance for each fixed rival beyond a rival-dependent
time implies the uniform theorem here. Fang and Chang's
\href{https://doi.org/10.1016/j.ipl.2014.06.012}{2014 paper} already
treats \(q=2\). Berend and Sher's
\href{https://doi.org/10.1080/15326349.2026.2671340}{2026 paper}
concerns Schilling's different question about uniformity over \emph{all}
packages. A targeted search for fixed-draw, eventual, and uniform
late-deadline results found no matching theorem; \textbf{this is not a
priority assessment or peer review}.

For fixed \(q,s\), only finitely many horizons remain between \(q+1\)
and \(256n\log(2n)\); this is not a finite reduction over all \(q,s\).
The bottleneck is the additive cross-cluster intersection loss \(R\),
unusable at short times where collisions occur frequently. The numerical
constant has not been optimized.

\end{document}
